\section{Background and Taxonomy}
\label{sec:background}

\subsection{Threat Model}

A guardrail operates within a threat model that specifies who the adversary is,
what access they have, and what harm they aim to cause.
We adopt the adversary taxonomy of Dong~et~al.~\cite{dong2024building} without
modification, as it provides adequate coverage for the real-time and multimodal
settings of this review and avoids introducing an incompatible competing
classification.
The three adversary classes are:

\begin{description}
  \item[Naive user] A user who elicits harmful outputs unintentionally or through
        naive prompting.
  \item[Jailbreaker] An adversary who deliberately crafts inputs to bypass safety
        constraints, with black-box API access only.
  \item[Fine-tuner] An adversary with the ability to customise the model (e.g.\ via a
        fine-tuning API), who degrades safety alignment during post-deployment
        customisation~\cite{qi2023finetuning}.
\end{description}

Real-time and multimodal settings expand the attack surface: a jailbreaker may
exploit streaming behaviour to extract partial harmful responses before a classifier
can fire, or embed harmful intent in images to evade text-focused filters.

\subsection{Guardrail Taxonomy}

We organise existing guardrail mechanisms along three axes:
\emph{placement in the inference pipeline}, \emph{implementation layer}, and
\emph{modality coverage}.
The axes are defined to be mutually orthogonal: placement describes \emph{when}
during a live inference request a guardrail fires; implementation layer describes
\emph{how} the guardrail is realised architecturally; and modality coverage
describes \emph{what input types} the guardrail can inspect.
Mechanisms whose safety behaviour is embedded in model weights during training
have \emph{no placement in the inference pipeline}---they impose no inference-time
interception point---and are therefore classified exclusively under the
implementation layer axis, not the placement axis.

\textit{Boundary case resolved}: plug-in adapter approaches such as
SafeVLM~\cite{oh2024safevlm} sit at the boundary between model-intrinsic and
model-extrinsic.
SafeVLM adds a safety projector, safety tokens, and a classification head that are
\emph{trained jointly} with or appended to the base VLM's architecture, which
blurs the weight-modification / external-classifier boundary.
We resolve this boundary by introducing a third implementation-layer subcategory,
\textbf{model-augmented}, defined as follows: a mechanism is model-augmented if and
only if (i) it adds new parameters to the base model's computational graph (rather
than wrapping the model in an external pipeline), \emph{and} (ii) the base model's
original weights are held fixed or can be restored by removing the added parameters.
This distinguishes model-augmented mechanisms from model-intrinsic mechanisms (which
modify the base weights) and from model-extrinsic mechanisms (which operate outside
the model's computational graph entirely, e.g., as a separate classifier invoked
after inference).
SafeVLM satisfies both criteria and is therefore classified as \textbf{model-augmented}
in this taxonomy.
MLLM-Protector~\cite{pi2024mllm} (Section~\ref{sec:multimodal}) similarly qualifies
as model-augmented under this criterion, as it adds adapter modules without altering
the host VLM's weights.
The implementation-layer axis thus has three values: \emph{model-intrinsic} (weights
modified), \emph{model-augmented} (parameters added, base weights preserved),
\emph{model-extrinsic} (external pipeline component).

\subsubsection{Placement in the Inference Pipeline}

This axis covers only \emph{inference-time} interception points, i.e.\ positions
at which a running guardrail can observe or interrupt a live inference request.
Training-time mechanisms (RLHF, constitutional AI, safety context distillation)
do not appear here because they impose no inference-time pipeline stage.

\begin{description}
  \item[Pre-generation (input)] The guardrail inspects the user prompt before any
        tokens are generated.
        Typical implementations are classifiers or rule engines applied to the prompt
        (e.g.\ Llama Guard in prompt-check mode~\cite{inan2023llama}).
  \item[In-generation (streaming)] The guardrail operates concurrently with token
        generation, either by inspecting each token or by inspecting growing
        prefixes.
        This is the least developed category and the central gap identified in
        this review.
  \item[Post-generation (output)] The guardrail inspects the completed response before
        it is returned to the user.
        Most LLM-as-judge and classifier-based systems operate here.
\end{description}

\noindent\textbf{Hybrid placement}: a mechanism may fire at more than one
pipeline stage.
NeMo Guardrails~\cite{rebedea2023nemo}, for instance, applies rails at both
the input stage (inspecting the user prompt) and the output stage (inspecting the
generated response), as well as at intermediate stages (knowledge retrieval,
action execution).
Such mechanisms are classified under \emph{all placement values that apply};
the taxonomy does not require a single placement assignment.
This multi-placement case should be interpreted as indicating that the mechanism
provides defence in depth across the pipeline, at the cost of additive latency
from each interception point.

\subsubsection{Implementation Layer}

\begin{sloppypar}
\begin{description}
  \item[Model-intrinsic (weight-embedded)] The mechanism modifies model weights or
        training procedure, baking safety behaviour into the model itself.
        No inference-time overhead is incurred, but the mechanism cannot be updated
        without retraining.
        Includes RLHF~\cite{ouyang2022training}, constitutional AI /
        RLAIF~\cite{bai2022constitutional}, and safety context distillation---the
        only weight-embedding technique with a fully reproducible description in the
        literature~\cite{touvron2023llama2}.
        Because these mechanisms have no inference pipeline placement, they do not
        appear in the placement axis above.
  \item[Model-augmented] The mechanism adds new parameters to the base model's
        computational graph but holds the base model's original weights fixed;
        the added modules can in principle be detached to restore the original
        model.
        This category is defined here to resolve boundary cases that neither
        model-intrinsic nor model-extrinsic classifies cleanly (see boundary
        case discussion above).
        Includes SafeVLM~\cite{oh2024safevlm} and
        MLLM-Protector~\cite{pi2024mllm}.
        Model-augmented mechanisms have an inference pipeline placement (typically
        pre- or post-generation, depending on where the added modules intercept).
  \item[Model-extrinsic] The mechanism is external to the model under guard and
        operates at inference time without adding parameters to the model's graph.
        Includes: (a) \emph{classifier-based} systems that run a separate safety
        classifier (Llama Guard~\cite{inan2023llama}, ShieldLM~\cite{zhang2024shieldlm},
        WildGuard~\cite{han2024wildguard}); (b) \emph{pipeline frameworks} that
        intercept at multiple stages (NeMo Guardrails~\cite{rebedea2023nemo}); and
        (c) \emph{rule engines} or prompt-injection defences.
        All model-extrinsic mechanisms have an inference pipeline placement (pre,
        in, or post-generation).
\end{description}
\end{sloppypar}

% SOTA Guardrail Taxonomy Diagram
% Three-axis taxonomy: Implementation Layer (vertical) x Placement (horizontal) x Modality (colour)
% Compile with: pdflatex main.tex (requires tikz, shapes.geometric, arrows.meta, positioning, calc)
%
% Grid coordinates (all in TikZ cm units):
%   Vertical dividers:   x = 1.1, 4.1, 7.1, 10.6, 13.2
%   Horizontal dividers: y = -2.0, -5.1, -7.5, -12.0
%   Column centres:      x = 2.6, 5.6, 8.85, 11.9
%   Row centres:         y = -3.55, -6.3, -9.75
%   Row label centre x:  -0.05

\begin{figure}[!htbp]
\centering
\resizebox{\linewidth}{!}{%
\begin{tikzpicture}[
  font=\small,
  textonly/.style={
    draw=blue!60, fill=blue!10, rounded corners=2pt,
    text width=2.0cm, align=center, minimum height=1.0cm,
    font=\scriptsize
  },
  multimodal/.style={
    draw=orange!70, fill=orange!12, rounded corners=2pt,
    text width=2.0cm, align=center, minimum height=1.0cm,
    font=\scriptsize
  },
  intrinsic/.style={
    draw=green!50!black, fill=green!10, rounded corners=2pt,
    text width=2.0cm, align=center, minimum height=1.0cm,
    font=\scriptsize
  },
  gap/.style={
    draw=red!60, fill=red!8, rounded corners=2pt, dashed,
    text width=2.0cm, align=center, minimum height=1.0cm,
    font=\scriptsize\itshape
  },
  arrow/.style={-Stealth, thick, black!60},
  divider/.style={dashed, black!30, thick},
  header/.style={fill=black!8, draw=black!70, thick},
]

%% -----------------------------------------------------------------------
%% AXIS ARROWS (drawn first so headers paint over them)
%% -----------------------------------------------------------------------
\draw[arrow, very thick] (-1.2, 0.5) -- (13.2, 0.5)
  node[midway, above, font=\footnotesize\bfseries]
  {Placement in inference pipeline $\longrightarrow$};

\draw[arrow, very thick] (-1.2, 0.5) -- (-1.2, -12.0)
  node[midway, left, font=\footnotesize\bfseries, rotate=90, yshift=6pt]
  {$\longleftarrow$ Implementation layer};

%% -----------------------------------------------------------------------
%% COLUMN HEADERS — TikZ rectangles at exact grid x-coords
%% x boundaries: 1.1 | 4.1 | 7.1 | 10.6 | 13.2
%% y range: 0.4 (below arrow) → -2.0 (header bottom = first h-divider top)
%% -----------------------------------------------------------------------
\fill[black!8] (1.1, 0.4) rectangle (13.2, -2.0);
% outer border
\draw[black!70, thick] (1.1,  0.4) -- (13.2,  0.4);
\draw[black!70, thick] (1.1, -2.0) -- (13.2, -2.0);
\draw[black!70, thick] (1.1,  0.4) -- (1.1,  -2.0);
\draw[black!70, thick] (13.2, 0.4) -- (13.2, -2.0);
% inner vertical borders (match grid dividers exactly)
\draw[black!70, thick] (4.1,  0.4) -- (4.1,  -2.0);
\draw[black!70, thick] (7.1,  0.4) -- (7.1,  -2.0);
\draw[black!70, thick] (10.6, 0.4) -- (10.6, -2.0);
% header text at column midpoints
\node[font=\footnotesize\bfseries, align=center] at (2.6,  -0.8) {Pre-generation\\(input)};
\node[font=\footnotesize\bfseries, align=center] at (5.6,  -0.8) {In-generation\\(streaming)};
\node[font=\footnotesize\bfseries, align=center] at (8.85, -0.8) {Post-generation\\(output)};
\node[font=\footnotesize\bfseries, align=center, text width=2.2cm] at (11.9, -0.8) {Training-time\\(no pipeline stage)};

%% -----------------------------------------------------------------------
%% ROW LABELS — TikZ rectangles at exact grid y-coords
%% x range: -1.2 → 1.1
%% y boundaries: -2.0 | -5.1 | -7.5 | -12.0
%% -----------------------------------------------------------------------
\fill[black!8] (-1.2, -2.0) rectangle (1.1, -12.0);
% outer border
\draw[black!70, thick] (-1.2, -2.0)  -- (1.1, -2.0);
\draw[black!70, thick] (-1.2, -12.0) -- (1.1, -12.0);
\draw[black!70, thick] (-1.2, -2.0)  -- (-1.2, -12.0);
\draw[black!70, thick] (1.1,  -2.0)  -- (1.1,  -12.0);
% inner horizontal borders (match grid dividers exactly)
\draw[black!70, thick] (-1.2, -5.1) -- (1.1, -5.1);
\draw[black!70, thick] (-1.2, -7.5) -- (1.1, -7.5);
% row label text at row midpoints  (x centre = (-1.2+1.1)/2 = -0.05)
\node[font=\footnotesize\bfseries, align=center] at (-0.05, -3.55) {Model\\Extrinsic};
\node[font=\footnotesize\bfseries, align=center] at (-0.05, -6.3)  {Model\\Augmented};
\node[font=\footnotesize\bfseries, align=center] at (-0.05, -9.75) {Model\\Intrinsic};

%% -----------------------------------------------------------------------
%% GRID LINES (drawn after fills so they appear on top)
%% -----------------------------------------------------------------------
% Horizontal dividers between content rows
\draw[divider] (1.1, -5.1) -- (13.2, -5.1);
\draw[divider] (1.1, -7.5) -- (13.2, -7.5);
% Vertical dividers between content columns
\draw[divider] (4.1,  -2.0) -- (4.1,  -12.0);
\draw[divider] (7.1,  -2.0) -- (7.1,  -12.0);
\draw[divider] (10.6, -2.0) -- (10.6, -12.0);
% Bottom border
\draw[divider] (-1.2, -12.0) -- (13.2, -12.0);

%% -----------------------------------------------------------------------
%% ROW 1: MODEL-EXTRINSIC  (y band: -2.0 to -5.1)
%% Pre-gen nodes at x=2.6; in-gen x=5.6; post-gen x=8.85
%% -----------------------------------------------------------------------
\node[textonly]  at (2.6,  -2.8) {Llama Guard\\\cite{inan2023llama}};
\node[textonly]  at (2.6,  -3.7) {ShieldLM\\\cite{zhang2024shieldlm}};
\node[textonly]  at (2.6,  -4.6) {NeMo Rails$^\dagger$\\\cite{rebedea2023nemo}};

\node[gap]       at (5.6,  -3.55) {Open gap:\\no streaming\\classifier};

\node[textonly]  at (8.85, -2.9) {WildGuard\\\cite{han2024wildguard}};
\node[textonly]  at (8.85, -4.1) {NeMo Rails$^\dagger$\\\cite{rebedea2023nemo}};

%% -----------------------------------------------------------------------
%% ROW 2: MODEL-AUGMENTED  (y band: -5.1 to -7.5, centre=-6.3)
%% -----------------------------------------------------------------------
\node[multimodal] at (2.6,  -6.3) {SafeVLM\\\cite{oh2024safevlm}};
\node[multimodal] at (8.85, -6.3) {MLLM-Protector\\\cite{pi2024mllm}};

%% -----------------------------------------------------------------------
%% ROW 3: MODEL-INTRINSIC  (y band: -7.5 to -12.0, training-time column x=11.9)
%% -----------------------------------------------------------------------
\node[intrinsic] at (11.9, -8.2)  {RLHF\\\cite{ouyang2022training}};
\node[intrinsic] at (11.9, -9.6)  {Const.\ AI\\\cite{bai2022constitutional}};
\node[intrinsic] at (11.9, -11.1) {Safety Context\\Distillation\\\cite{touvron2023llama2}};

%% -----------------------------------------------------------------------
%% LEGEND (modality)
%% -----------------------------------------------------------------------
\node[textonly,   text width=1.5cm, minimum height=0.5cm] at (2.6,  -12.8) {Text-only};
\node[multimodal, text width=1.5cm, minimum height=0.5cm] at (5.6,  -12.8) {Multimodal};
\node[intrinsic,  text width=1.5cm, minimum height=0.5cm] at (8.85, -12.8) {Weight-embedded};
\node[gap,        text width=1.5cm, minimum height=0.5cm] at (11.9, -12.8) {Open gap};
\node[font=\footnotesize\bfseries] at (-0.05, -12.8) {Modality:};

\end{tikzpicture}%
}% end resizebox
\caption{%
  Three-axis taxonomy of LLM guardrail mechanisms.
  \textbf{Horizontal axis}: placement in the inference pipeline
  (pre-generation, in-generation/streaming, post-generation, training-time).
  \textbf{Vertical axis}: implementation layer
  (model-extrinsic, model-augmented, model-intrinsic).
  \textbf{Colour}: modality coverage
  (\textcolor{blue!70}{text-only}, \textcolor{orange!80!black}{multimodal}, \textcolor{green!50!black}{weight-embedded}).
  The dashed red cell marks the central open gap identified in this review:
  no streaming-compatible extrinsic guardrail currently exists.
  $^\dagger$NeMo Guardrails~\cite{rebedea2023nemo} applies rails at both
  pre-generation (input inspection) and post-generation (output inspection);
  it appears in both columns to reflect this hybrid placement.
  Its latency cost grows additively with each interception point.
}
\label{fig:sota_taxonomy}
\end{figure}


\subsubsection{Modality Coverage}

Table~\ref{tab:modality} summarises the modality coverage of the principal systems
and approaches reviewed.

\begin{table}[htbp]
  \caption{Modality coverage of nine representative guardrail systems.
           Systems selected to span all three taxonomy axes (placement,
           implementation layer, and modality); selection is not exhaustive.
           DRESS excluded: not a dedicated guardrail (see text; appears in
           Table~\ref{tab:vlm_defence}).
           \checkmark = addressed; $\triangle$ = partial/unverified;
           -- = not addressed.}
  \label{tab:modality}
  \centering
  \begin{tabular}{lcccc}
    \toprule
    \textbf{System / Approach} & \textbf{Text} & \textbf{Image} & \textbf{Audio} & \textbf{Video} \\
    \midrule
    RLHF~\cite{ouyang2022training}        & \checkmark & -- & -- & -- \\
    Constitutional AI~\cite{bai2022constitutional} & \checkmark & -- & -- & -- \\
    Safety Context Distillation~\cite{touvron2023llama2} & \checkmark & -- & -- & -- \\
    Llama Guard~\cite{inan2023llama}      & \checkmark & $\triangle$ & -- & -- \\
    NeMo Guardrails~\cite{rebedea2023nemo} & \checkmark & -- & -- & -- \\
    ShieldLM~\cite{zhang2024shieldlm}    & \checkmark & -- & -- & -- \\
    WildGuard~\cite{han2024wildguard}    & \checkmark & -- & -- & -- \\
    SafeVLM~\cite{oh2024safevlm}          & \checkmark & $\triangle$ & -- & -- \\
    Azure Content Safety~\cite{azure2024contentsafety} & \checkmark & \checkmark & -- & -- \\
    \bottomrule
  \end{tabular}
\end{table}

The table reveals that audio and video safety is entirely absent from the verified
academic literature, and that image safety is addressed only partially or without
reproducible evaluation.
DRESS~\cite{chen2024dresssafety} is excluded from this table because it is not a
dedicated VLM safety guardrail; it appears in Section~\ref{sec:multimodal} as a
training-time alignment baseline contrast, where its peripheral status is documented
in full.

\subsection{Taxonomy Cell Analysis}

A three-axis taxonomy with values \{pre-generation, in-generation, post-generation\} $\times$ \{model-intrinsic, model-augmented, model-extrinsic\} $\times$ \{text, image, audio, video\} implies up to $3 \times 3 \times 4 = 36$ cells.
Not all cells are structurally realizable.

\paragraph{Structurally impossible cells}
Model-intrinsic mechanisms (weight-embedded) operate during training, not during inference; the placement axis is undefined for them.
The 12 placement-specific cells in the model-intrinsic row are training-time / placement-agnostic: model-intrinsic mechanisms (RLHF, constitutional AI, safety context distillation) operate during training, not during inference-time placement, and are therefore agnostic to the pre/in/post-generation placement axis.
These collapse to 4 placement-agnostic cells (one per modality), reducing the total to 28 cells: 24 placement-specific (model-augmented and model-extrinsic rows) and 4 placement-agnostic (model-intrinsic row).

Two boundary cases require explicit treatment: system prompt engineering, though often classified as model-intrinsic, operates at inference time and is more accurately placed in the pre-generation $\times$ model-extrinsic cell; activation steering, a weight-space intervention applied during the forward pass, is genuinely model-intrinsic and in-generation, and is accordingly treated as placement-agnostic in the model-intrinsic row.
(Here, activation steering refers to the post-training weight-space or representation-space modification, not to inference-time hidden-state intervention; inference-time activation steering, which modifies hidden states during the forward pass, would be classified as in-generation $\times$ model-intrinsic and is noted as a boundary case not yet represented in the corpus.)

\paragraph{Populated cells (from the 75-paper corpus)}
Of the 24 placement-specific cells, the following are populated by at least one corpus paper:
pre-generation $\times$ model-extrinsic $\times$ text (Llama Guard, ShieldLM, WildGuard, AEGIS);
post-generation $\times$ model-extrinsic $\times$ text (same systems in response-checking mode);
pre-generation $\times$ model-augmented $\times$ image (SafeVLM~\cite{oh2024safevlm});
post-generation $\times$ model-augmented $\times$ text+image (MLLM-Protector~\cite{pi2024mllm});
pre+post-generation $\times$ model-extrinsic $\times$ text (NeMo Guardrails, hybrid placement).
Systems listed in multiple placement cells (e.g., Llama Guard, which supports both pre- and post-generation deployment) occupy a single row in the populated-cell count; placement multiplicity reflects deployment flexibility, not multiple distinct mechanisms.

\paragraph{Structurally empty by research absence}
Two voids are systematic and constitute the central findings of this review.
Within the 28-cell framework, \emph{the in-generation row} (streaming) is entirely unpopulated across the placement-specific cells: no paper in the corpus evaluates or proposes a model-augmented or model-extrinsic mechanism that operates on the growing token prefix during active generation (Section~\ref{sec:realtime}, Streaming Safety Gap).
\emph{The audio and video modality columns} are unpopulated across all 28 cells regardless of placement or implementation layer: no paper addresses audio-language or video-language safety mechanisms at any pipeline stage (Section~\ref{sec:multimodal}, Audio and Video Safety).
Together, these two voids --- the in-generation row and the audio/video columns --- define the research agenda developed in Section~\ref{sec:open_problems}.

\subsection{Foundational Alignment Techniques}

Three intrinsic alignment techniques form the background against which extrinsic
guardrails are evaluated.

\paragraph{RLHF}
Reinforcement learning from human feedback~\cite{ouyang2022training} fine-tunes an
LLM using a reward model trained on human preference labels.
The technique produced the first commercially deployed aligned assistants but
requires large-scale human annotation and is susceptible to reward hacking.

\paragraph{Constitutional AI}
Bai~et~al.~\cite{bai2022constitutional} replace human safety labels with a set of
written principles (the ``constitution'') from which an AI generates and filters
its own training data.
Of the 18 quantitative and architectural claims extracted from the background and
taxonomy sources (see Table~\ref{tab:claim_counts}), two survived structured
self-review: the Safety Context Distillation mechanism and the dual reward model
arbitration threshold documented in Section~\ref{sec:realtime}.
The claim that no human labels are required did not survive structured self-review
on three specific grounds: (i) the ``constitution'' itself is a
human-authored document requiring substantive expert judgement about which principles
to include; (ii) the RL from AI Feedback (RLAIF) stage uses a preference model that
is initialised from human feedback data; and (iii) Bai~et~al.\ report using human
evaluation to calibrate and validate the resulting model's harmlessness, contradicting
the no-human-labels framing.
The technique is a genuine advance in reducing annotation volume, but it does not
eliminate human involvement from the alignment process.

\paragraph{Safety Context Distillation}
Introduced in Llama~2~\cite{touvron2023llama2}, this technique generates safe
responses by prefixing prompts with a safety preprompt, fine-tunes the model on
those responses, then removes the preprompt.
The result is a model that has internalised safety behaviour in its weights.
This is the only weight-embedding technique whose mechanism is described with
sufficient precision to be reproducible; its use of the safety reward model to
select which samples receive distillation is a concrete, verifiable architectural
decision.
